% =====================================================================
%  2. Gradient and directional derivative
% =====================================================================
\section{The gradient: which way is downhill?}

\begin{frame}{Gradient and directional derivative}
  \begin{itemize}\small
    \item The gradient collects all partial derivatives: $\nabla f(\x) = \bigl(\tfrac{\partial f}{\partial x_1},\dots,\tfrac{\partial f}{\partial x_n}\bigr)\T$.
    \item The directional derivative along $\vd$ tells you how fast $f$ changes if you move that way
          (chain rule on $t\mapsto f(\x+t\vd)$):
  \end{itemize}
  \[
    D_{\vd} f(\x) = \lim_{t\to0}\frac{f(\x + t\vd) - f(\x)}{t} = \nabla f(\x)\T\vd .
  \]
  \begin{itemize}\small
    \item First-order Taylor expansion: $f(\x + t\vd) = f(\x) + t\,\nabla f(\x)\T\vd + o(t)$.
    \item So if $\nabla f(\x)\T\vd < 0$, a small step along $\vd$ lowers $f$. We call $\vd$ a \alert{descent direction}.
  \end{itemize}
  \begin{keybox}[The question that starts everything]
    Of all the unit directions you could step in, which one lowers $f$ the \emph{fastest}?
  \end{keybox}
\end{frame}

\begin{frame}{Theorem: $-\nabla f$ is the steepest way down}
  \begin{columns}[T]
    \begin{column}{0.55\textwidth}
      \begin{thmbox}[Theorem (steepest descent)]
        If $\nabla f(\x)\neq\mathbf 0$, then
        \[
          \min_{\norm{\vd}=1}\nabla f(\x)\T\vd = -\norm{\nabla f(\x)},
        \]
        and the only minimiser is $\vd^\star = -\nabla f(\x)/\norm{\nabla f(\x)}$.
      \end{thmbox}
      \small\textbf{Proof.}
      \pause
      By Cauchy--Schwarz, $|\nabla f\T\vd| \le \norm{\nabla f}\norm{\vd} = \norm{\nabla f}$, so
      $\nabla f\T\vd \ge -\norm{\nabla f}$.
      \pause
      Cauchy--Schwarz is an equality only when $\vd$ is parallel to $\nabla f$. Taking the minus sign,
      $\vd = -\nabla f/\norm{\nabla f}$ reaches the bound. $\blacksquare$
    \end{column}
    \begin{column}{0.43\textwidth}
      \centering
      % Steepest-descent geometry: the half-plane of descent directions, the unit
% circle of candidate directions d, the gradient and the angle theta.
\begin{tikzpicture}[scale=1.4, >={Stealth[length=7pt]}]
  % descent half-plane  {d : grad^T d < 0}  (gradient direction (0.447, 0.894))
  \begin{scope}
    \clip (-1.75,-1.35) rectangle (1.75,1.35);
    \fill[mGreen!9] (-2,1) -- (2,-1) -- (2,-3) -- (-2,-3) -- cycle;
  \end{scope}
  \draw[mGrey, dashed, line width=0.6pt] (-1.75,0.875) -- (1.75,-0.875);
  \node[font=\scriptsize, mGreen, anchor=south east] at (1.60,-1.23) {descent directions};
  \node[font=\scriptsize, mGrey, rotate=-26.6, anchor=south west] at (-1.72,0.94) {level set};
  % unit circle
  \draw[mDarkTeal!35, line width=1pt] (0,0) circle (1);
  % candidate directions (faint)
  \foreach \t in {200,230,260,290,320,350,80,110,140,170}
    \draw[mGrey!45, -{Stealth[length=4pt]}] (0,0) -- (\t:1);
  % a generic direction d and the angle
  \draw[->, line width=1.3pt, mBlue] (0,0) -- (-15:1) node[right, font=\small] {$\vd$};
  \draw[mBlue, line width=0.7pt] (63.4:0.43) arc[start angle=63.4, end angle=-15, radius=0.43];
  \node[mBlue, font=\small] at (24:0.67) {$\theta$};
  % gradient and steepest descent
  \draw[->, line width=1.8pt, mRed] (0,0) -- (63.4:1.18) node[above, font=\small] {$\nabla f(\x)$};
  \draw[->, line width=1.8pt, mGreen] (0,0) -- (243.4:1) node[below left=3pt, font=\small] {$\vd^\star$};
  \fill[mInk] (0,0) circle (1.4pt);
  \node[font=\small, mInk, fill=mPaper, inner sep=1pt] at (-0.28,0.27) {$\x$};
\end{tikzpicture}


      \footnotesize $D_{\vd}f = \norm{\nabla f}\cos\theta$, which is smallest at $\theta = \pi$.
      Every arrow in the green half-plane goes downhill; $\vd^\star$ goes down fastest.
    \end{column}
  \end{columns}
\end{frame}

\begin{frame}{The gradient is perpendicular to the contours}
  \begin{columns}[T]
    \begin{column}{0.47\textwidth}
      \small
      Take any smooth curve $\gamma(t)$ that stays on the level set $\{f = c\}$, with $\gamma(0) = \x$. Then
      \[
        f(\gamma(t)) = c \;\Rightarrow\; \nabla f(\x)\T\gamma'(0) = 0 .
      \]
      So $\nabla f(\x)$ is perpendicular to every direction along the contour.

      \vspace{4pt}
      \begin{itemize}\footnotesize
        \item The arrows show $-\nabla f$ (normalised) for $f = x^4+y^4-4xy$.
        \item Look at how they cross every contour at a right angle, run away from the saddle along $(1,1)$,
              and pour into the two minima.
      \end{itemize}
    \end{column}
    \begin{column}{0.51\textwidth}
      \centering
      \includegraphics[width=0.92\linewidth]{fig_landscape_contour}
    \end{column}
  \end{columns}
\end{frame}

\begin{frame}{``Steepest'' depends on how you measure length}
  Swap the Euclidean norm for $\norm{\vd}_{\bP} = \sqrt{\vd\T\bP\vd}$ with $\bP\succ0$, and the answer changes:
  \[
    \min_{\vd\T\bP\vd\le1}\nabla f\T\vd
    \quad\Longrightarrow\quad
    \vd^\star \propto -\bP^{-1}\nabla f(\x).
  \]
  \small\textbf{Proof.} Substitute $\vd = \bP^{-1/2}\mathbf z$ with $\norm{\mathbf z}\le1$. The objective becomes
  $(\bP^{-1/2}\nabla f)\T\mathbf z$, so the previous theorem gives $\mathbf z^\star\propto-\bP^{-1/2}\nabla f$ and therefore
  $\vd^\star\propto-\bP^{-1}\nabla f$. $\blacksquare$

  \vspace{4pt}
  \begin{center}
  \begin{tabular}{@{}l l@{}}
    \toprule
    Choice of $\bP$ & What you get \\
    \midrule
    $\bI$ & plain gradient descent \\
    a diagonal scaling & feature standardisation, Adam-type methods (Section~13) \\
    $\nabla^2 f(\x)$ & the Newton direction (a different topic, mentioned only) \\
    \bottomrule
  \end{tabular}
  \end{center}
  \footnotesize\textcolor{mGrey}{So there is no single ``steepest'' direction. It depends on the geometry you pick~\cite{boyd2004},
  and choosing that geometry well is what preconditioning does (Section~8).}
\end{frame}
